\section{Validation}
\label{chap:proposal:validation}

The validation of \acs{5G} network applications cannot be limited to the testing methodologies of conventional software systems. This evaluation must verify the application functions as intended and effectively utilizes the exposed \acs{5G} capabilities, just deployed as services. However, these aspects cannot be meaningfully assessed in isolation, since the behavior of a network application results from the interplay of multiple factors: (i) its internal logic, (ii) how it interacts with the network through the available exposure interfaces, and (iii) ultimately the network conditions and capabilities that are actually delivered by the \acs{5G} system as a result of these interactions. 

In real operational environments, these conditions are highly dynamic and variable, often deviating from what's expected by a given vertical application, due to factors such as traffic load and the shared physical infrastructure. Network applications must therefore be capable of adapting to such variability, which must also be accounted for during their validation. This includes testing the application under both optimal (high throughput, low latency and jitter, no packet loss, etc.) and degraded scenarios (congestion, jitter, packet loss, loss of connectivity, etc.). Only through such comprehensive testing is it possible to assess the robustness and performance of the application and to determine whether it can satisfy the defined \acsp{SLA} when deployed over a real \acs{5G} network.

To achieve this, two different validation approaches are suggested: (i) simulating network events at the network exposure layer to create controlled and reproducible test scenarios; or (ii) relying on a real \acs{5G} network to introduce network impairments, directly affecting the behavior of the application under test.

The first approach focuses on the validation of the application logic, and how the application should react to the events and responses returned by the \acs{5G} network. By emulating responses such as analytics values, event notifications, or the outcome of \acs{QoS} requests, a wide range of scenarios can be evaluated without having to reproduce them in the actual network. In fact, this approach does not even require the existence of a \acs{5GC} or a fully operational \acs{5G} network. Since the \acs{NEF} is the entry point to \acs{5G} network programmability, it can be adapted to provide such responses through its standardized \acsp{API}. Open-source solutions, such as NEFSim~\cite{NEFsim}, have been developed specifically for this purpose. This is especially valuable at early stages of development, to assess the functional validity of the application before it is exercised on the actual network. In addition, should this approach be adopted in the testbed, its benefits would extend to other available exposure capabilities that internally consume the \acs{NEF}, such as CAMARA \acsp{API}. Therefore, the same emulation model could be used to validate the interaction with the network, independently of the exposure model preferred by the application.

The second approach evaluates the application's behavior under real network conditions provided by a deployed \acs{5G} system, once again taking advantage of the flexibility and programmability offered by the testbed. Through the network exposure interfaces described earlier, experimenters can request the enforcement of specific network conditions, i.e. \acs{QoS} profiles, and establish different networking scenarios with distinct performance characteristics. Under each scenario, the application is expected to adapt its behavior to the conditions enforced by the network: under favorable profiles, it should fully exploit the available resources and meet its target \acsp{KPI}, whereas under degraded conditions it should react gracefully, for instance, by lowering its service level (e.g. streaming resolution) or requesting a higher-priority \acs{QoS} profile. 

Moreover, the testbed capabilities could also be used to deliberately induce adverse or changing network conditions the application is expected to withstand. By stress testing the network with competing traffic, through its active probing mechanism, the testbed can reproduce these conditions in the live network, effectively degrading the \acs{QoS} of the provisioned slice and reducing the bandwidth available to the application, without relying on changes to its \acs{QoS} profile.

This section outlines the plan for the validation of the vertical use case presented in Section~\ref{chap:proposal:use-case}, in the proposed testbed. This includes the complete flow required to validate the application, including the service provisioning process needed to order and deploy the application and the testbed's services, which is detailed first. Furthermore, by validating the network application and exercising the testbed's network exposure interfaces and capabilities under realistic conditions, the testbed itself will also be validated, demonstrating its ability to support the end-to-end deployment and operation of vertical applications.

\subsection{Service Delivery}
\label{chap:proposal:validation:delivery}

Service delivery is the first step of the experimentation process, ordering and provisioning the underlying environment for the evaluation of vertical applications. For the \acs{PoC}, this environment would be composed of multiple services offered by the catalog presented in Section~\ref{chap:proposal:delivery:design}. The experimenter would first onboard the network application, packaging the application backend, the inference server and the \acs{PPDR} operator dashboards as a Network Application service.

This service would declare dependencies on the experiment's remaining services, parameterized to the use case. The Network Slice service would configure a \acs{NEST} meeting the uplink bitrate and latency requirements of \ref{enum:nfr:uplink-bitrate}. The User Equipment service would register the camera \acs{UE}, attaching it to the provisioned slice. The Network Exposure service would include \acs{CAPIF}, the \acs{NEF}, and optionally CAMARA \acsp{API}, according to the chosen abstraction level. Finally, the Network Observability service would provision the observability dashboards and, optionally, a prober to actively monitor the camera. At this point, the validation should confirm that the characteristics exposed by the catalog services suffice to express the requirements of the experimenter, according to the service design proposed in Section~\ref{chap:proposal:delivery:design}.

The experiment would then be instantiated through a single service order, submitted through the \acs{GUI} of the \acs{OSS} or the \acs{TMF}~641 \acs{API}. The validation of this process consists of verifying that the \acs{OSS} provisions each service following the chain depicted in Figure~\ref{proposal:delivery:orchestration-lifecycle}, deploying each phase only once its prerequisites are active. Upon completion, the state of the service order and the service inventory, retrieved through the \acs{TMF}~638 and \acs{TMF}~641 \acsp{API}, respectively, should indicate that all configured services are active and that their relationships are properly registered. Similarly, the resource inventory (\acs{TMF}~639) should confirm the provisioning of the underlying resources, matching the parameters of the original order.

To ensure that the resources have been successfully provisioned, it is also necessary to verify that the delegated orchestrator has deployed the corresponding resources. For instance, for containerized services orchestrated, for example, by Kubernetes, this would involve confirming, through its native \acs{API}, that the expected pods were scheduled and that their containers are running and ready. This lower level resources should also reflect the attributes of the ordered service, e.g. those of the ordered \acs{NEST}. Nonetheless, the conformance of the network conditions delivered by the slice should be assessed separately, as described in Section~\ref{chap:proposal:validation:use-case}.

A failure in any phase should prevent the provisioning of subsequent services, and the termination of the experiment should remove all of them. Since all application components would be deployed through the \acs{OSS}, without manual intervention, validating these aspects also verifies~\ref{enum:nfr:packaging}. Furthermore, repeatedly instantiating the same services across multiple executions should yield equivalent service and resource inventories, demonstrating consistent provisioning and proving the reproducibility of the service delivery process.

\subsection{Multi-Vertical Use Case}
\label{chap:proposal:validation:use-case}

Once the required services have been provisioned, the envisioned network application for the multi-vertical use case will be exercised and validated end-to-end on the testbed. In particular, this validation seeks to evaluate how the testbed enables the application to request and enforce the capabilities and network conditions demanded by the use case, and whether it can correctly adapt its behavior as these conditions change.

As noted at the beginning of this section, the functioning of a network application results from the combined operation of three complementary components, all of which must be considered during validation: (i) the application itself, (ii) the exposure layer, formed by the \acs{NEF} and CAMARA services, and (iii) the underlying \acs{5G} network. The testbed should therefore provide experimenters with the evidence required to understand how these components interact in pursuit of the application's goals. This requires analyzing observability data (e.g. metrics, logs) collected at each point of the experimental chain and correlating them, so that the observed behavior of the application can be traced back to its causes. This allows experimenters to verify not only that each request issued by the application was accepted, but also that it translated into the expected network conditions, and that the application reacted to those conditions as intended. For example, when an experimenter orders a network slice, the provisioned slice is expected to deliver the throughput and latency expressed in the attributes of the ordered \acs{NEST}. Verifying these conditions involves comparing the requested \acs{QoS} values with the metrics collected for the experiment, including those obtained through active probing.
